% Macro per la struttura dei template dei casi d'uso.
% Richiede: tabularx, enumitem (gia' caricati nel preambolo principale).
%
% NOTA: tabularx (colonne di tipo X) scansiona i token "&" e "\\" senza
% espanderli, per calcolare la larghezza delle celle. Per questo motivo
% \begin{tabularx}...\end{tabularx} e il "\\ \hline" di fine riga NON
% possono essere nascosti dentro una macro o un environment: vanno
% scritti letteralmente in ciascun file che usa queste macro. Le macro
% coprono solo il contenuto della singola riga (etichetta in grassetto
% + "&" + contenuto).

\newcommand{\ucstart}{\renewcommand{\arraystretch}{1.5}\noindent}

\newcommand{\ucid}[1]{\textbf{ID Caso d'Uso} & \textbf{#1}}
\newcommand{\ucname}[1]{\textbf{Nome} & #1}
\newcommand{\ucactor}[1]{\textbf{Attore Primario} & #1}
\newcommand{\ucsecondaryactor}[1]{\textbf{Attore Secondario} & #1}
\newcommand{\ucdescription}[1]{\textbf{Descrizione} & #1}
\newcommand{\ucprecond}[1]{\textbf{Pre-condizioni} & #1}
\newcommand{\ucpostcond}[1]{\textbf{Post-condizioni} & #1}
\newcommand{\ucmainflow}[1]{%
  \textbf{Flusso Principale} &
  \begin{enumerate}[topsep=0pt,partopsep=6pt,before=\vspace*{-\baselineskip}]
    #1
  \end{enumerate}%
}
\newcommand{\ucaltflow}[1]{\textbf{Flussi Alternativi} & #1}
\newcommand{\ucnote}[1]{\textbf{Note} & #1}
\newcommand{\ucvariations}[1]{\textbf{Variazioni} & #1}
